\subsection{Applications: II. Extension, to a space of measures, of a continuous function with values in a space of operators}

Let G be a Hausdorff locally convex space, H a Hausdorff and quasi-complete locally convex space, and denote by F the space \(\mathcal{L}(\mathrm{G};\mathrm{H})\) of continuous linear mappings of G into H, equipped with the topology of compact convergence. The space F is not necessarily quasi-complete, and if \(t \mapsto U(t)\) is a continuous mapping of T into F, and \(\mu\) is a positive measure on T with compact support, one does not necessarily have \(\int U d\mu \in F\) (Exer. 27). However, if, for every compact subset K of T, \(U(\mathrm{K})\) is equicontinuous, then its balanced convex envelope in F is also equicontinuous (TVS, III, §3, No.~4), and since H is quasi-complete, the closure of this balanced convex envelope will be a complete subset of F (TVS, III, §3, No.~8, Prop. 11); one will then indeed have \(\int U d\mu \in F\) (No.~2, Prop. 8).

The supplementary condition imposed on U may be expressed otherwise:

Lemma 3. --- Let G, H be two locally convex spaces, U a mapping of a locally compact space T into \(\mathcal{L}(\mathrm{G};\mathrm{H})\) . The following conditions are equivalent:

\begin{enumerate}
\def\labelenumi{\alph{enumi})}
\item
  The mapping \((t,\mathbf{x})\mapsto U(t)\cdot \mathbf{x}\) of \(\mathrm{T}\times \mathrm{G}\) into H is continuous.
\item
  For every compact subset K of T, \(U(\mathbf{K})\) is equicontinuous, and there exists a total set \(D \subset G\) such that for every \(x \in D\) , the mapping \(t \mapsto U(t) \cdot x\) is continuous on T.
\end{enumerate}

Moreover, when \(U\) satisfies these conditions, \(U\) is a continuous mapping of \(\mathrm{T}\) into \(\mathcal{L}(\mathrm{G};\mathrm{H})\) equipped with the topology of compact convergence.

To see that a) implies b), we observe that for every neighborhood V of 0 in H and every \(t \in \mathbf{K}\) , there exist by hypothesis a neighborhood \(\mathrm{L}_t\) of \(t\) in T and a neighborhood \(\mathrm{W}_t\) of 0 in G such that the relations \(t' \in \mathrm{L}_t\) and \(\mathbf{x} \in \mathrm{W}_t\) imply \(U(t') \cdot \mathbf{x} \in \mathrm{V}\) . It suffices to cover K with a finite number of neighborhoods \(\mathrm{L}_{t_i}\) and to take \(\mathrm{W} = \bigcap_{i} \mathrm{W}_{t_i}\) in order to have \(U(t) \cdot \mathbf{x} \in \mathrm{V}\) whenever \(t \in \mathrm{K}\) and \(\mathbf{x} \in \mathrm{W}\) , which proves the equicontinuity of \(U(\mathrm{K})\) .

Conversely, suppose b) is verified; it suffices to show that for every compact subset K of T, the mapping \((t,\mathbf{x})\mapsto U(t)\cdot\mathbf{x}\) is continuous on \(K\times G\) . Let \(M=U(K)\) ; since M is equicontinuous, it follows that on M, the topology of pointwise convergence in G is identical to the topology of pointwise convergence in D (GT, X, §2, No.~4, Th. 1); the hypothesis b) therefore implies that \(t\mapsto U(t)\) is a continuous mapping of K into \(\mathcal{L}(G;H)\) when \(\mathcal{L}(G;H)\) is equipped with the topology of pointwise convergence. On the other hand, \((A,\mathbf{x})\mapsto A\cdot\mathbf{x}\) is a continuous mapping of \(M\times G\) into H when M is equipped with the topology of pointwise convergence (GT, X, §2, No.~1, Cor. 4 of Prop. 1). Since the mapping \((t,\mathbf{x})\mapsto U(t)\cdot\mathbf{x}\) may be factored as \((t,\mathbf{x})\mapsto(U(t),\mathbf{x})\mapsto U(t)\cdot\mathbf{x}\) , we conclude that it is continuous.

Finally, the last assertion of the lemma follows from the fact that, on M, the topology of compact convergence is identical to that of pointwise convergence (GT, X, §2, No.~4, Th. 1).

Thus, suppose that U satisfies the conditions of Lemma 3; then (if H is quasi-complete) one defines, as in No.~6, a linear mapping \(\lambda \mapsto \int U d\lambda\) of \(\mathcal{C}'(\mathrm{T})\) into \(\mathrm{F} = \mathcal{L}(\mathrm{G}; \mathrm{H})\) . We set \(U(\lambda) = \int U d\lambda\) .

PROPOSITION 16. --- Let G, H be two Hausdorff locally convex spaces, with H assumed to be quasi-complete. Let U be a mapping of T into \(\mathcal{L}(\mathrm{G};\mathrm{H})\) such that \((t,\mathbf{x})\mapsto U(t)\cdot \mathbf{x}\) is a continuous mapping of \(\mathrm{T}\times \mathrm{G}\) into H. Then the bilinear mapping \((\lambda ,\mathbf{x})\mapsto U(\lambda)\cdot \mathbf{x}\) of \(\mathcal{C}'(\mathrm{T})\times \mathrm{G}\) into H is hypocontinuous relative to the equicontinuous subsets of \(\mathcal{C}'(\mathrm{T})\) and the compact subsets of G (which implies that the linear mapping \(\lambda \mapsto U(\lambda)\) of \(\mathcal{C}'(\mathrm{T})\) into F is continuous).

The continuity of \(\lambda \mapsto U(t)\) as a mapping of \(\mathcal{C}'(\mathrm{T})\) into \(\mathrm{F}\) follows from Lemma 3 and Remark 2 following Prop. 14 of No.~6. Thus, it remains to prove that for every closed, balanced, convex neighborhood \(\mathrm{V}\) of 0 in \(\mathrm{H}\) and for every equicontinuous subset \(\mathrm{N}\) of \(\mathcal{C}'(\mathrm{T})\) , there exists a neighborhood \(\mathrm{W}\) of 0 in \(\mathrm{G}\) such that the relations \(\mathbf{x} \in \mathrm{W}\) , \(\lambda \in \mathrm{N}\) imply that \(U(\lambda) \cdot \mathbf{x} \in \mathrm{V}\) . One can suppose that \(\mathrm{N} = \mathrm{S}^{\circ}\) , where \(\mathrm{S}\) is a neighborhood of 0 in \(\mathcal{C}(\mathrm{T})\) , consequently one can suppose that \(\mathrm{S}\) is the set of functions \(g \in \mathcal{C}(\mathrm{T})\) such that \(|g(t)| \leqslant 1\) on a compact subset \(\mathrm{K}\) of \(\mathrm{T}\) . It suffices to show that \(|\langle U(\lambda) \cdot \mathbf{x}, \mathbf{x}' \rangle| \leqslant 1\) for \(\mathbf{x} \in W\) , \(\mathbf{x}' \in V^\circ\) and \(\lambda \in S^\circ\) . Now, since \(U(K)\) is equicontinuous, there exists a neighborhood \(W\) of 0 in \(G\) such that the relations \(t \in K\) , \(\mathbf{x} \in W\) imply \(U(t) \cdot \mathbf{x} \in V\) ; the relations \(\mathbf{x} \in W\) , \(\mathbf{x}' \in V^\circ\) therefore imply that the function \(t \mapsto \langle U(t) \cdot \mathbf{x}, \mathbf{x}' \rangle\) belongs to \(S\) , hence that \(|\langle U(t) \cdot \mathbf{x}, \mathbf{x}' \rangle| = |\int\langle U(t) \cdot \mathbf{x}, \mathbf{x}' \rangle d\lambda(t)| \leqslant 1\) by the definition of \(S^\circ\) .

Let us now assume that \(U\) is a continuous mapping of \(\mathrm{T}\) into \(\mathrm{F}\) and, moreover, that \(U(\mathrm{T})\) is equicontinuous. Then, the same reasoning as above shows (since \(\mathrm{H}\) is quasi-complete) that for every bounded positive measure \(\mu\) on \(\mathrm{T}\) , \(\int U d\mu \in \mathrm{F}\) . One can therefore define, as above, a linear mapping \(\lambda \mapsto \int U d\lambda = U(\lambda)\) of \(\mathcal{M}^1(\mathrm{T})\) into \(\mathrm{F}\) that extends the analogous mapping of \(\mathcal{C}'(\mathrm{T})\) into \(\mathrm{F}\) . Moreover, for every closed, balanced, convex neighborhood \(\mathrm{V}\) of 0 in \(\mathrm{H}\) , there exists by hypothesis a neighborhood \(\mathrm{W}\) of 0 in \(\mathrm{G}\) such that for every \(\mathbf{x} \in \mathrm{W}\) and every \(t \in \mathrm{T}\) , one has \(U(t) \cdot \mathbf{x} \in \mathrm{V}\) , consequently (since \(\mathrm{V}\) is weakly closed) \(\int (U(t) \cdot \mathbf{x}) d\lambda(t) \in \| \lambda \| \cdot \mathrm{V}\) (No.~2, Prop. 5). In other words:

PROPOSITION 17. --- Let G, H be two Hausdorff locally convex spaces, with H assumed to be quasi-complete. Let U be a mapping of T into \(\mathcal{L}(G;H)\) such that \((t,\mathbf{x})\mapsto U(t)\cdot \mathbf{x}\) is continuous on T×G, and \(U(T)\) is equicontinuous. Then, if \(\mathcal{M}^1 (T)\) is equipped with its Banach space topology, the bilinear mapping \((\lambda ,\mathbf{x})\mapsto U(\lambda)\cdot \mathbf{x}\) of \(\mathcal{M}^1 (T)\times G\) into H is continuous (which implies, in particular, that the linear mapping \(\lambda \mapsto U(\lambda)\) of \(\mathcal{M}^1 (T)\) into \(\mathcal{L}(G;H)\) is continuous when \(\mathcal{L}(G;H)\) is equipped with the topology of bounded convergence).

PROPOSITION 18. --- Let \(G_{1}, G_{2}, H_{1}, H_{2}\) be four Hausdorff locally convex spaces, with \(H_{1}\) and \(H_{2}\) assumed to be quasi-complete. Let \(A : G_{1} \to G_{2}\) and \(B : H_{1} \to H_{2}\) be two continuous linear mappings. Let \(U_{1} : T \to \mathcal{L}(G_{1}; H_{1})\) , \(U_{2} : T \to \mathcal{L}(G_{2}; H_{2})\) be two mappings satisfying the conditions of Prop. 16 (resp. Prop. 17), and suppose that for every \(t \in T\) , \(B \circ U_{1}(t) = U_{2}(t) \circ A\) . Then, for every measure with compact support (resp. every bounded measure) \(\lambda\) on T, \(B \circ U_{1}(\lambda) = U_{2}(\lambda) \circ A\) .

Indeed, for every \(x \in G_{1}\) , one has (No.~1, Prop. 1)

\[
\begin{array}{c} \left(B \circ U _ {1} (\lambda)\right) \cdot \mathbf {x} = \int \Big (\left(B \circ U _ {1} (t)\right) \cdot \mathbf {x} \Big) d \lambda (t) \\ = \int \Big (\left(U _ {2} (t) \circ A\right) \cdot \mathbf {x} \Big) d \lambda (t) = U _ {2} (\lambda) \cdot (A \cdot \mathbf {x}). \end{array}
\]

Remarks. --- 1) Suppose that G and H are Banach spaces, and let U be a mapping of T into \(\mathcal{L}(\mathrm{G};\mathrm{H})\) such that \((t,\mathbf{x})\mapsto U(t)\cdot\mathbf{x}\) is continuous on \(T\times G\) . Note that this implies that the finite function \(t\mapsto\|U(t)\|\) is bounded on every compact subset of T and lower semi-continuous on T, being the upper envelope of the continuous functions \(t \mapsto |U(t) \cdot \mathbf{x}|\) as x runs over the ball \(|x| \leqslant 1\) in G. Set \(h(t) = \|U(t)\|\) . Then, for every positive measure \(\mu\) on T such that h is \(\mu\) -integrable, we again have \(\int U d\mu \in \mathcal{L}(\mathrm{G};\mathrm{H})\) . For, the measure \(\nu = h \cdot \mu\) is bounded by hypothesis; there therefore exists a partition of T formed by a \(\nu\) -negligible set N and a sequence \((\mathrm{K}_{n})\) of compact subsets. The argument made at the beginning of this No., applied to the measure \(\varphi_{K_{n}} \cdot \mu\) , shows that

\[
A _ {n} = \int \varphi_ {\mathrm{K} _ {n}} U d \mu \in \mathrm{F} = \mathcal {L} (\mathrm{G}; \mathrm{H}),
\]

and, moreover (No.~2, Prop. 6), \(\| A_n\| \leqslant \int \varphi_{\mathrm{K}_n}\| U\| d\mu \leqslant \nu (\mathrm{K}_n)\) . The series with general term \(A_{n}\) is therefore absolutely convergent in the Banach space \(\mathcal{L}(\mathrm{G};\mathrm{H})\) , and it is immediate that its sum is \(\int Ud\mu\) and that \(\left\| \int Ud\mu \right\| \leqslant \int \| U\| d\mu\) .

\begin{enumerate}
\def\labelenumi{\arabic{enumi})}
\setcounter{enumi}{1}
\tightlist
\item
  Suppose that \(\mathrm{G} = \mathrm{H}\) is quasi-complete and that \(U\) satisfies the hypotheses of Prop. 16. Let \(\mathbf{M}\) be a dense subset of the space \(\mathcal{C}'(\mathrm{T})\) , for the weak topology \(\sigma\big(\mathcal{C}'(\mathrm{T}),\mathcal{C}(\mathrm{T})\big)\) , and let \(\mathrm{X}\) be a closed linear subspace of \(\mathrm{H}\) such that \(U(\lambda)(\mathrm{X})\subset \mathrm{X}\) for every measure \(\lambda \in \mathrm{M}\) . Then also \(U(t)(\mathrm{X})\subset \mathrm{X}\) for every \(t\in \mathrm{T}\) : indeed, for every \(\mathbf{x}\in \mathrm{X}\) and every \(\mathbf{x}'\in \mathrm{H}'\) orthogonal to \(\mathrm{X}\) , by hypothesis \(\langle U(\lambda)\cdot \mathbf{x},\mathbf{x}'\rangle = 0\) for all \(\lambda \in \mathrm{M}\) , which may be written \(\int \langle U(t)\cdot \mathbf{x},\mathbf{x}'\rangle d\mu (t) = 0\) . The continuous function \(t\mapsto \langle U(t)\cdot \mathbf{x},\mathbf{x}'\rangle\) , being orthogonal to \(\mathrm{M}\) , is therefore 0, which yields \(\langle U(t)\cdot \mathbf{x},\mathbf{x}'\rangle = 0\) for every \(\mathbf{x}'\in \mathrm{X}^\circ\) , whence \(U(t)\cdot \mathbf{x}\in \mathrm{X}\) for all \(t\in \mathrm{T}\) and \(\mathbf{x}\in \mathrm{X}\) , and this proves our assertion.
\end{enumerate}

